\documentclass{article}
\usepackage{amssymb,amsmath,amsfonts,amsthm,setspace,graphicx,mathtools,tikz-cd,tikz,xcolor}
\usepackage[hyphens]{url}
\usepackage{hyperref}

\newtheoremstyle{noperiod}
  {}
  {}
  {\normalfont}
  {}
  {\bfseries}
  { }
  {.5em}
  {}

\theoremstyle{noperiod}
\newtheorem{definition}{Definition}
\newtheorem{proposition}{Proposition}
\newtheorem{lemma}{Lemma}
\newtheorem{theorem}{Theorem}
\newtheorem{example}{Example}
\newtheorem{corollary}{Corollary}
\newtheorem{remark}{Remark}

\makeatletter
\renewenvironment{proof}[1][\proofname]{\par
  \pushQED{\qed}
  \normalfont \topsep6\p@\@plus6\p@\relax
  \trivlist
  \item[\hskip\labelsep
        \itshape
    #1 ]\ignorespaces
}{
  \popQED\endtrivlist\@endpefalse
}
\makeatother

\newcommand{\notmid}{\mathrel{\ooalign{$\mkern-5mu\not$\crcr$|$}}}

\renewcommand{\proofname}{\bfseries Proof}
\renewcommand{\qedsymbol}{$\blacksquare$}

\title{CanonicalForms}
\date{}
\author{}
\begin{document}
\maketitle

\begin{abstract}
We formulate, in the language of set theory, the concepts of invariant and
canonical form for an equivalence relation $\sim$ on a set $X$.
A map $f \colon X \to Y$ is called an invariant of $\sim$ if $x \sim y
\implies f(x) = f(y)$, and a complete invariant if equality of images
also implies the equivalence relation.
A canonical form is an invariant $C \colon X \to X$ satisfying $x \sim
C(x)$ for all $x \in X$; its image $C(X)$ is called the skeleton or
complete system of representatives of $\sim$.
We characterise canonical forms via retraction--section pairs, showing that $C$
is a canonical form if and only if there exists a section $s$ of the quotient
map $q \colon X \to X/\sim$ such that $C = s \circ q$.
An analogous characterisation holds for surjective complete invariants.
We prove that every complete invariant $I \colon X \to Y$ factors uniquely
through the surjective complete invariant $I' \colon X \to I(X)$, and define
$s' \circ I'$ as the canonical form of $\sim$ determined by the complete
invariant $I$, whenever a section $s'$ of $I'$ exists.
Assuming the axiom of choice, we show that the canonical form determined by any
complete invariant exists.
Propositions that assume the axiom of choice are marked
with~$^{\textrm{AC}}$.
\end{abstract}

\begin{definition}
Let $X, Y$ be sets. Let $f \colon X \to Y$ be a map and let $\sim$ be an
equivalence relation on $X$. We say that $f$ is an \textbf{invariant} of $\sim$
if the following condition holds.
\begin{center}
  ${^{\forall}x,y} \in X,\ (x \sim y \implies f(x) = f(y))$
\end{center}
We further say that $f$ is a \textbf{complete invariant} of $\sim$ if the
following condition holds.
\begin{center}
  ${^{\forall}x,y} \in X,\ (x \sim y \iff f(x) = f(y))$
\end{center}
\hfill $\square$
\end{definition}

\begin{example}
Let $X$ be a set. Let $\sim$ be an equivalence relation on $X$ and let
$q \colon X \to X/\sim$ be the quotient map. Then $q$ is a complete invariant
of $\sim$.
\hfill $\square$
\end{example}

\begin{definition}
Let $X, Y$ be sets. Let $C \colon X \to X$ be an invariant of an equivalence
relation $\sim$ on $X$. We say that $C$ is a \textbf{canonical form} of $\sim$
if the following condition holds.
\begin{center}
  ${^{\forall}x} \in X,\ x \sim C(x)$
\end{center}
We call the image $C(X)$ of $X$ under $C$ the \textbf{skeleton} or the
\textbf{complete system of representatives} of $\sim$.
\hfill $\square$
\end{definition}

\begin{proposition}
Let $X, Y$ be sets. Let $C \colon X \to X$ be a canonical form of an
equivalence relation $\sim$ on $X$. Then $C = C \circ C$.
\hfill $\square$
\end{proposition}

\begin{proof}
Let $x \in X$. Since $x \sim C(x)$, we have $C(x) = C(C(x))$. Therefore
$C = C \circ C$.
\end{proof}

\begin{definition}
Let $X, Y$ be sets and let $r \colon X \to Y$ and $s \colon Y \to X$ be maps.
We call the pair $(r, s)$ a \textbf{retraction--section pair} if the following
condition holds. We call $r$ the \textbf{retraction} of $s$, and $s$ the
\textbf{section} of $r$.
\begin{center}
  $r \circ s = \mathrm{id}_Y$
\end{center}
\hfill $\square$
\end{definition}

\begin{lemma}
Let $X$ be a set and let $q \colon X \to X/\sim$ be the quotient map of an
equivalence relation $\sim$ on $X$. For a map $C \colon X \to X$, the following
conditions \textup{(1)} and \textup{(2)} are equivalent.
\def\labelenumi{(\theenumi)}
\begin{enumerate}
  \item $C$ is a canonical form of $\sim$.
  \item There exists a section $s \colon X/\sim \to X$ of $q \colon X \to X/\sim$
        such that $C = s \circ q$.
\end{enumerate}
\hfill $\square$
\end{lemma}

\begin{proof}
$(1) \Rightarrow (2)$:\; Assume $C$ is a canonical form of $\sim$. For each
$y \in X/\sim$, since $q$ is surjective there exists $x \in X$ such that
$y = q(x)$. Set $s(y) = C(x)$. For any $x' \in X$ with $q(x) = q(x')$, we
have $x \sim x'$, so $C(x) = C(x')$. Thus $s \colon X/\sim \to X$ is
well-defined. For any $x \in X$, $C(x) = s(q(x))$, so $C = s \circ q$. For
any $x \in X$, $(q \circ s)(q(x)) = q(s(q(x))) = q(C(x))$, and since
$x \sim C(x)$ we have $q(C(x)) = q(x)$. Therefore $q \circ s =
\mathrm{id}_{X/\sim}$.

$(2) \Rightarrow (1)$:\; Assume there exists a section $s \colon X/\sim \to X$
of $q \colon X \to X/\sim$ such that $C = s \circ q$. For $x, x' \in X$, if
$x \sim x'$, then $C(x) = (s \circ q)(x) = s(q(x)) = s(q(x')) = (s \circ q)(x')
= C(x')$, so $C$ is an invariant of $\sim$. For any $x \in X$,
$q(x) = \mathrm{id}_{X/\sim}(q(x)) = (q \circ s)(q(x)) = q((s \circ q)(x)) =
q(C(x))$, so $x \sim C(x)$. Therefore $C$ is a canonical form of $\sim$.
\end{proof}

\begin{proposition}
Let $X, Y$ be sets and let $I \colon X \to Y$ be a surjective complete invariant
of an equivalence relation $\sim$ on $X$. For a map $C \colon X \to X$, the
following conditions \textup{(1)} and \textup{(2)} are equivalent.
\def\labelenumi{(\theenumi)}
\begin{enumerate}
  \item $C$ is a canonical form of $\sim$.
  \item There exists a section $s' \colon Y \to X$ of $I \colon X \to Y$ such
        that $C = s' \circ I$.\hfill $\square$
\end{enumerate}
\end{proposition}

\begin{proof}
By Lemma~1, it suffices to show that the existence of a section
$s \colon X/\sim \to X$ of $q \colon X \to X/\sim$ with $C = s \circ q$ is
equivalent to~(2). Define $\bar{I} \colon X/\sim \to Y$ by $\bar{I}(q(x)) = I(x)$
for all $x \in X$. When $x \sim x'$, since $I$ is an invariant we have
$I(x) = I(x')$, so $\bar{I}$ is well-defined. For any $x, x' \in X$, if
$I(x) = I(x')$ then, since $I$ is a complete invariant, $x \sim x'$, hence
$q(x) = q(x')$, so $\bar{I}$ is injective. For any $y \in Y$, since
$I \colon X \to Y$ is surjective there exists $x \in X$ with $I(x) = y$, giving
$\bar{I}(q(x)) = y$. Therefore $\bar{I}$ is bijective and hence invertible.

Suppose there exists a section $s$ of $q$ with $C = s \circ q$. Set
$s' = s \circ \bar{I}^{-1}$. Since $q = \bar{I}^{-1} \circ I$, we have
$C = s \circ \bar{I}^{-1} \circ I = s' \circ I$. Moreover,
$I \circ s' = (\bar{I} \circ q) \circ (s \circ \bar{I}^{-1})
= \bar{I} \circ (q \circ s) \circ \bar{I}^{-1}
= \bar{I} \circ \mathrm{id}_{X/\sim} \circ \bar{I}^{-1}
= \mathrm{id}_Y$,
so $s'$ is a section of $I$.

Conversely, suppose~(2) holds. Set $s = s' \circ \bar{I}$. Then
$s \circ q = s' \circ \bar{I} \circ q = s' \circ I = C$.
Since $I \circ s' = \mathrm{id}_Y$, we have $I \circ s' \circ I = I$,
and therefore $I \circ C = I$.
Since $I = \bar{I} \circ q$, it follows that $\bar{I} \circ q \circ C = \bar{I} \circ q$,
and by injectivity of $\bar{I}$ we obtain $q \circ C = q$.
For any $y \in X/\sim$ there exists $x \in X$ such that $y = q(x)$.
Then
\[
(q \circ s)(y)
= (q \circ s' \circ \bar{I})(y)
= (q \circ s' \circ \bar{I} \circ q)(x)
= (q \circ s' \circ I)(x)
= (q \circ C)(x)
= q(x) = y,
\]
so $q \circ s = \mathrm{id}_{X/\sim}$.
\end{proof}

\begin{proposition}
Let $X, Y$ be sets and let $I \colon X \to Y$ be a complete invariant of an
equivalence relation $\sim$ on $X$. The unique map $I' \colon X \to I(X)$
satisfying $I = i \circ I'$ with respect to the inclusion $i \colon I(X) \to Y$
is surjective and is a complete invariant of $\sim$.
\hfill $\square$
\end{proposition}

\begin{proof}
Define $I' \colon X \to I(X)$ by $I'(x) = I(x)$ for all $x \in X$. Since
$I = i \circ I'$ and $I'(X) = I(X)$, $I'$ is surjective. For any $x, x' \in X$,
$x \sim x' \iff I(x) = I(x') \iff I'(x) = I'(x')$,
so $I'$ is a complete invariant of $\sim$. If $I'' \colon X \to I(X)$ satisfies
the condition, then $I'(x) = I(x) = I''(x)$ for all $x \in X$, so $I'$ is
unique.
\end{proof}

\begin{definition}
Let $X, Y$ be sets and let $I \colon X \to Y$ be a complete invariant of an
equivalence relation $\sim$ on $X$. When a section $s' \colon I(X) \to X$ of
the unique map $I' \colon X \to I(X)$ satisfying $I = i \circ I'$ with respect
to the inclusion $i \colon I(X) \to Y$ exists, we call the canonical form
$s' \circ I'$ of $\sim$ the \textbf{canonical form of $\sim$ determined by the
complete invariant $I$}.
\hfill $\square$
\end{definition}

\begin{lemma}\hspace{-0.7em}$^{\textrm{AC}}$\hspace{0.2em}
Let $X, Y$ be sets and let $f \colon X \to Y$ be a surjection. Then there
exists a map $g \colon Y \to X$ such that $f \circ g = \mathrm{id}_Y$.
\hfill $\square$
\end{lemma}

\begin{proof}
Since $f \colon X \to Y$ is surjective, $f^{-1}(y) \neq \emptyset$ for every
$y \in Y$. Therefore, by the axiom of choice, $\prod_{y \in Y} f^{-1}(y) \neq
\emptyset$, so there exists $g \in \prod_{y \in Y} f^{-1}(y)$. This map
$g \colon Y \to X$ satisfies $f \circ g = \mathrm{id}_Y$.
\end{proof}

\begin{proposition}\hspace{-0.7em}$^{\textrm{AC}}$\hspace{0.2em}
Let $X, Y$ be sets and let $\sim$ be an equivalence relation on $X$. Then the
canonical form of $\sim$ determined by the complete invariant $I$ exists.
\hfill $\square$
\end{proposition}

\begin{proof}
Let $q \colon X \to X/\sim$ be the quotient map. Since $q$ is surjective, by
Lemma~2 there exists a map $s \colon X/\sim \to X$ satisfying $q \circ s =
\mathrm{id}_{X/\sim}$. By Lemma~1, $s \circ q$ is a canonical form of $\sim$,
and by Proposition~2, taking $I'$ to be the surjection determined by $I$, there
exists a section $s' \colon I(X) \to X$ of $I'$ such that $s \circ q = s' \circ
I'$. Then $s' \circ I'$ is the canonical form of $\sim$ determined by the
complete invariant $I$.
\end{proof}

\end{document}